\chapter*{Introduction}
\addcontentsline{toc}{chapter}{Introduction}


A \emph{proof assistant} is a software tool that provides a formal language for
defining mathematical objects, stating their properties, and constructing
proofs to verify those claims. The system then rigorously verifies
these proofs down to their fundamental logical axioms.

While traditionally used to verify software correctness, proof assistants
are increasingly applied to abstract mathematics—a primary focus of the
mathlib community. Through formalization, every definition becomes strictly
unambiguous, ensuring that all resulting proofs are virtually guaranteed
to be free of errors. The Lean theorem prover is a proof assistant
developed principally by Leonardo de Moura. The latest version is called Lean 4.
More of the technical details can be found in \cite{deMoura2021}.

Built on this foundation is \emph{mathlib}, a highly active, community-driven
initiative aiming to create a comprehensive, unified library of formalized
mathematics in Lean. Alongside its mathematical core, the library includes
foundational definitions for programming, sustained by daily updates from a
dedicated network of contributors.

Since the
establishment of the Lean Focused Research Organization \cite{fro2024} in 2023, the Lean 4 ecosystem has experienced explosive growth. This
has made Lean an ideal tool not only for students undertaking smaller
mathematical projects but also the platform of choice for scientists managing
large-scale collaborative endeavors. Early landmark achievements, such as
\emph{Liquid Tensor Experiment} \cite{commelin2022liquid} and
there is a related work about the categorical foundations of formalizated of
condensed mathematics in \cite{asgeirsson2024categorical}.
% Moreover, the
% formalization of sphere eversions in differential topology \cite{vandoorn2023},
% demonstrated the feasibility of handling complex, abstract mathematics in Lean.

Between 2024 and 2026, formalized mathematics witnessed several milestone
breakthroughs. In the field of analysis, a team led by Floris van Doorn
completed the complete formalization of Carleson's theorem and its
generalizations, overcoming long-standing computational
complexity bottlenecks in fourier analysis. This is a large cooperation project and its blue print can be found in
\cite{becker2025blueprintformalizationcarlesonstheorem}. In number theory, Fermat's Last
Theorem for regular primes is completely formalized in 2025 and their paper can be
found in \cite{best2025completeformalizationfermatstheorem}.

Simultaneously, the intersection of auto-formalization and artificial
intelligence has emerged as the frontier of growth in this field. An
epoch-making development was the verification of the sphere packing problem
in late 2025. Building upon Maryna Viazovska's work
\cite{viazovska2017sphere},
researchers utilized large language models to assist in code translation,
rapidly completing a ``sorry-free'' formalized verification of
optimal sphere packings
in dimensions 8 and 24 \cite{hariharan2026milestoneformalizationspherepacking}.

Given the results on the formalization of substitution and evaluation of
multivariate power series in \cite{chambertloir2025formalizationdividedpowerslean},
this thesis aims to formalize the definitions and applications of formal
group laws in Lean. A standard reference for formal group laws
is \cite{hazewinkel1978formal}. Our mathematical discussion will mainly
be based on \cite{hazewinkel1978formal, cassels1967algebraic, LubinTate1965,
silverman2009arithmetic}.

In Chapter 1, we will discuss the theory of formal group laws from a purely
mathematical perspective. In Chapter 2, we will discuss Lubin--Tate theory,
which is one of the most important applications of formal group laws. Chapter
3 provides a brief introduction to the aspects of Lean necessary for
understanding the formalization of Chapter 1's content, which we will
then present in Chapter 4. The code can be found at
https://github.com/WenrongZou/FormalGroupLaws.
Throughout this thesis, we will link to code from either our formalization
or Mathlib. These links will be marked with this symbol: \href{https://github.com/WenrongZou/FormalGroupLaws/blob/main/FormalGroupLaws/Basic.lean}{\faExternalLink}.

% Given the results of formalization of subtitution and evaluation
% of multivariate power series
% in \cite{chambertloir2025formalizationdividedpowerslean}.
% This thesis aims to
% formalize definitions and applications formal group laws in Lean. A standard reference
% of formal group laws is \cite{hazewinkel1978formal}. Our mathematical discussion will
% mainly be based on \cite{hazewinkel1978formal}, \cite{cassels1967algebraic},
% \cite{LubinTate1965} and \cite{silverman2009arithmetic}.

% In chapter 1, we will discuss the theory of formal group laws from a purely
% mathematical perspective. In chapter 2, we will discuss Lubin--Tate theory, which is
% one of the most important applications of formal group laws. Chapter 3 gives a brief
% introduction to some aspects of Lean that will be useful to understand the
% formalisation of somem of content of chapter 1 which we will cover in chapter 4.
% The accompanying code can be found at
% https://github.com/WenrongZou/FormalGroupLaws. Throughout this thesis we will link
% to code either from our formalisation or from mathlib. These links will be mark with this
% symbol: \href{https://github.com/WenrongZou/FormalGroupLaws/blob/main/FormalGroupLaws/Basic.lean}{\faExternalLink}.